\clearpage
\thispagestyle{empty}
\pdfbookmark[0]{Vorschriften zur Bestimmung der magnetischen Wirkung, die ein Magnetstab in der Ferne ausübt. Ausgabe und Quellen}{magnetic-effect-at-distance-1840-sources}
\begin{center}\Large\textbf{Vorschriften zur Bestimmung der magnetischen Wirkung, die ein Magnetstab in der Ferne ausübt}\end{center}

Carl Friedrich Gauss’ Abhandlung erschien in den \textit{Resultaten aus den Beobachtungen des magnetischen Vereins} für 1840 als Artikel II. Die KI-Übertragung der Einleitung und aller sieben Artikel in modernes Deutsch mit KI-erstelltem Satz folgt dem deutschen Text in den \textit{Werken}, Bd.~V (Göttingen, 1867), S.~427--435. Die Koordinaten, Vorzeichen, historischen Einheiten und beide Zahlentafeln folgen dieser Ausgabe.

\medskip\noindent
Zur Terminologie siehe die Einführung in den Erdmagnetismus des U.S. Geological Survey und die Übersicht zum Erdmagnetfeld des GFZ Helmholtz-Zentrums für Geoforschung. Für die Koordinaten gelten in dieser Abhandlung die von Gauss verwendeten Konventionen.

\medskip\noindent
\href{https://gdz.sub.uni-goettingen.de/id/PPN236006339}{Gesammelte deutsche Ausgabe (1867)}\par
\noindent\href{https://doi.org/10.1007/978-3-642-49319-5_20}{Deutscher Nachdruck (1877)}\par
\noindent\href{https://www.usgs.gov/programs/geomagnetism/introduction-geomagnetism}{U.S. Geological Survey, Introduction to Geomagnetism}\par
\noindent\href{https://www.gfz.de/sektion/geomagnetismus/daten-produkte-dienste/haeufig-gestellte-fragen-zum-thema-erdmagnetismus}{GFZ, Häufig gestellte Fragen zum Thema Erdmagnetismus}
\clearpage
